\chapter{方法、建模或关键规则}

\section{本章作用}

前一章给出了 Pebble 的核心思想与设计原则，本章进一步回答“系统如何把这些原则落成可执行的方法结构”。这里的“方法”主要指对象如何被建模、输入如何被解释、结果如何沉淀，以及不同功能模块之间如何组成最小闭环。它既要足够抽象，能够解释全局设计；也要足够贴近系统对外呈现的输入、输出与数据边界，避免停留在口号层。

因此，本章不把 Pebble 写成一个单独的模型调用接口，而是把它写成一个由关系对象、当前输入、分析结果、训练记录与长期洞察共同构成的方法系统。后续流程章与系统章会分别展开执行路径与工程结构；本章先固定其方法对象、作用关系与关键规则。

\section{对象建模：Pebble 实际处理的五类核心对象}

Pebble 当前的方法结构可以压缩为五类核心对象。它们对应的并非抽象概念，而是当前工作区中已经可观察到的服务输入、输出与持久化载体。

\begin{table}[H]
  \centering
  \caption{Pebble 方法层的核心对象及其在系统中的作用。}
  \label{tab:pebble-method-objects}
  \begin{tabularx}{\textwidth}{>{\raggedright\arraybackslash}p{2.5cm} >{\raggedright\arraybackslash}X >{\raggedright\arraybackslash}X}
    \toprule
    对象 & 在方法层中的定义 & 当前系统中的典型承载位置 \\
    \midrule
    关系对象 & 被持续维护的上下文单元，包含关系类型、对方特点、期望结果、情境补充与生成上下文 & 关系服务与关系节点数据 \\
    当前输入 & 用户此刻需要理解或回应的文本、消息或对话轮次 & 解码请求文本、模拟陪练中的当前回合消息 \\
    分析结果 & 系统对当前输入形成的结构化解释，包括潜台词、风险类型、情绪得分与回应候选 & 解码服务输出、陪练侧栏反馈 \\
    练习记录 & 用户从分析或模拟中选取、保存并准备回看的应对内容 & 练习本条目与收藏结果 \\
    长期洞察 & 围绕具体关系逐步累积的观察摘要、生成上下文与策略适配条件 & 关系档案中的结构化字段与生成上下文 \\
    \bottomrule
  \end{tabularx}
\end{table}

这五类对象之间的关系决定了 Pebble 的方法边界。

\begin{itemize}
  \item 关系对象负责提供情境解释条件。
  \item 当前输入负责触发一次局部分析或一轮练习互动。
  \item 分析结果负责把原始输入转换成用户可理解、可比较、可操作的结构。
  \item 练习记录负责把一次性的建议转化为可积累的经验材料。
  \item 长期洞察负责把多轮互动中的局部观察重新汇回关系对象，形成下一次分析的先验条件。
\end{itemize}

\section{最小方法闭环}

在对象建模明确之后，Pebble 的方法可以被概括为一条四步闭环：条件注入、局部分析、反馈转化、长期沉淀。图~\ref{fig:pebble-method-loop} 对这一闭环做了压缩表达。

\begin{figure}[H]
  \centering
  \begin{tikzpicture}[
    node distance=1.7cm and 1.1cm,
    every node/.style={font=\small},
    methodbox/.style={
      draw=ReportBlue,
      rounded corners=10pt,
      fill=ReportSoftBlue,
      minimum width=3.1cm,
      minimum height=1.2cm,
      align=center
    },
    methodarrow/.style={
      -{Latex[length=3mm]},
      thick,
      draw=ReportBlue
    }
  ]
    \node[methodbox] (relation) {关系对象\\上下文条件};
    \node[methodbox, right=of relation] (decode) {当前输入\\解码或陪练分析};
    \node[methodbox, right=of decode] (feedback) {结构化反馈\\建议与评分};
    \node[methodbox, below=of decode] (practice) {练习记录\\保存与回看};

    \draw[methodarrow] (relation) -- (decode);
    \draw[methodarrow] (decode) -- (feedback);
    \draw[methodarrow] (feedback) -- (practice);
    \draw[methodarrow] (practice.west) |- (relation.south);
    \draw[methodarrow] (feedback.south) -- (practice.north);
  \end{tikzpicture}
  \caption{Pebble 的最小方法闭环：关系上下文进入当前分析，分析结果再回流为练习材料与长期洞察条件。}
  \label{fig:pebble-method-loop}
\end{figure}

这条闭环中的每一步都有明确的方法含义。

\begin{itemize}
  \item 条件注入：系统先收集与当前关系相关的最小必要上下文，使分析不只依赖孤立文本。
  \item 局部分析：系统对当前输入形成结构化解释，产出潜台词、场景判断、情绪得分与可选回应。
  \item 反馈转化：系统把分析结果转化为可练习、可选择、可保存的反馈对象，而不是停留在一次性展示。
  \item 长期沉淀：系统把练习与关系信息重新汇回关系对象，使后续分析逐步获得更稳定的上下文支撑。
\end{itemize}

\section{关键规则}

上述闭环若要真正成立，需要至少四条关键规则维持一致性。

\subsection{规则一：所有局部分析都允许显式上下文化}

当前解码服务的输入结构已经体现出这一规则：除了当前文本本身，方法层还保留了自由上下文与关系对象标识两个扩展入口。其意义在于，系统默认允许“当前文本 + 关系信息”的联合解释，而不是把上下文视为后处理补充。

这一规则带来两层约束。

\begin{itemize}
  \item 分析入口必须允许关系对象进入主路径，否则“关系上下文优先”只能停留在文档叙述层。
  \item 关系上下文应以最小必要信息组织，既能帮助解释，又不把无关信息无限扩展为提示噪声。
\end{itemize}

\subsection{规则二：输出必须结构化，而非仅返回自由文本}

Pebble 的方法要求分析结果至少包含三类信息：对当前话语的解释、对风险或情绪状态的压缩判断、以及可供用户实际使用的回应候选。陪练模块同样遵循这一逻辑，即使它表现为对话界面，仍然持续产出结构化评分、即时反馈与注意点。

结构化输出的重要性在于，它决定了结果能否进入下一步。

\begin{itemize}
  \item 若结果只有一段泛化说明，用户难以比较不同策略，也难以把某条回应保存为练习材料。
  \item 若结果具备固定槽位，例如潜台词、得分、方案与反馈，那么它就可以稳定进入练习本、统计与后续复盘路径。
\end{itemize}

\subsection{规则三：练习材料是方法闭环中的中间对象}

当前练习服务把解码结果与模拟陪练结果收敛为同一种练习本对象，这说明 Pebble 的方法并不把“保存”视作附属功能，而是把它视为从局部理解走向长期能力积累的关键中间环节。

在方法层，练习材料至少承担三种作用。

\begin{itemize}
  \item 它把当前一次分析中的候选回应固定下来，避免结果随着页面关闭而消失。
  \item 它为用户提供跨场景比较与复盘的基础，使不同来源的应对策略可以进入同一经验空间。
  \item 它为后续长期洞察提供间接证据，帮助系统理解哪些策略在什么关系对象和场景下更常被采纳。
\end{itemize}

\subsection{规则四：长期洞察以描述性上下文累积，而非一次性定性}

关系服务当前维护的字段结构说明，Pebble 的长期建模中心是“关系对象”，而非对人格或状态做最终判定。关系类型、对方特点、期望结果、情境补充和生成上下文共同构成一个可被持续修订的解释条件集合。该集合在方法上的价值，在于它允许系统把局部分析结果重新放回关系语境中理解。

这一规则直接约束了长期建模的表达方式。

\begin{itemize}
  \item 长期洞察应优先使用描述性信息和可回看的观察线索。
  \item 生成上下文应服务于后续解释与陪练，不应被包装成权威性诊断结论。
  \item 关系对象的更新应允许用户修订和重写，以保持用户主权与边界克制。
\end{itemize}

\section{方法层的综合理解}

综合来看，Pebble 的方法并非由某一个算法名称定义，而是由一组稳定的对象关系与规则关系定义。它把关系对象作为长期条件，把当前输入作为局部触发，把结构化分析作为中介，把练习记录作为沉淀载体，并把长期洞察作为下一轮分析的先验。这个方法结构能够同时解释解码、陪练、练习本与关系档案为何需要共同存在，也为下一章展开算法与执行路径提供了统一入口。
